\documentclass[10pt,a4paper]{article}
\usepackage[margin=1.6cm]{geometry}
\usepackage{xcolor}
\usepackage{hyperref}
\usepackage{amsmath,amssymb}
\usepackage{enumitem}
\usepackage[breakable]{tcolorbox}
\usepackage{booktabs}
\usepackage{fontspec}
\usepackage[CJKmath=true]{xeCJK}
\setCJKmainfont{Noto Sans CJK SC}[BoldFont=Noto Sans CJK SC Bold, AutoFakeBold=false]

\definecolor{cblue}{HTML}{4C72B0}
\definecolor{cgold}{HTML}{DD8452}
\definecolor{cgreen}{HTML}{55A868}
\definecolor{cred}{HTML}{C44E52}
\definecolor{bgskill}{HTML}{F0F7FF}

\hypersetup{colorlinks=true, linkcolor=cblue, urlcolor=cblue}
\setlength{\parskip}{4pt}\setlength{\parindent}{0pt}

\begin{document}

\begin{tcolorbox}[colback=bgskill, colframe=cblue, title=\textbf{NoveltySkill — Idea Card (L3 Oral-Ready)}, breakable]

\textbf{Title}: Finite-Sample Sharpening of the SFT-as-Adversarial-Regularizer Equivalence

\textbf{Hook}: \emph{ The SFT-loss-as-adversarial-regularizer equivalence holds only in the limit. We sharpen it to finite-sample form --- and the residual term is exactly the overoptimization signal RLHF training should regulate. }

\section*{Abstract}
Reward-model overoptimization in RLHF is the canonical Goodhart's-law failure: the policy KL constraint to the SFT prior holds, the proxy reward keeps climbing, but a held-out gold reward plateaus then degrades. Current mitigations --- constrained-RL thresholds (dblp:conf/iclr/MoskovitzSSSSDM24), iterative preference smoothing, ensemble-disagreement penalties --- are heuristic because the underlying theoretical equivalence between the SFT loss and an adversarial reward regularizer (openalex:W4399115631) holds only asymptotically (n $\to$ ∞), while production preference datasets (HH-RLHF n $\approx$ 161k pairs against million-parameter policy gradients) sit far from this limit. We sharpen the asymptotic equivalence into a finite-sample form by replacing the limit step with a Bernstein-type concentration bound on the preference-gradient covariance $\Sigma$\_pref, yielding a closed-form residual $\epsilon$(n, $\|$$\nabla$R\_proxy − $\nabla$R\_true$\|$\_$\Sigma$\_pref) computable from preference data alone. We propose using $\epsilon$ as the dynamic regularization schedule $\beta$\_t for SFT-anchored PPO. Falsification: train Llama-3-8B with the four schedules (vanilla, Moskovitz, ensemble, ours) on HH-RLHF using Llama-3-70B as gold proxy; the schedule fails kill-switch if (a) peak gold-reward improvement < 0.05 bits/token vs baselines or (b) empirical $\epsilon$ exceeds theoretical bound by > 3$\times$ at any of three checkpoints. Expected: $\ge$ 0.15 bits/token improvement and $\le$ 1.5$\times$ tracking at all checkpoints.

\section*{Motivation}
\textbf{Problem framing.} Reinforcement Learning from Human Feedback (RLHF) trains a policy against a learned reward model that imperfectly approximates the true reward. The classical Goodhart's-law signature emerges: the proxy reward grows, the policy's KL distance to the SFT prior stays bounded, but a held-out gold reward (from a stronger frozen reference model) plateaus and degrades. The user's RLHF pipeline exhibits this canonical pattern. Current literature presents three families of mitigation. Constrained RLHF (Moskovitz et al., ICLR 2024) reformulates as a constrained RL problem with Lagrangian thresholds on the proxy gap. Iterative Data Smoothing (Zhu et al., ICML 2024) attacks the upstream preference-data overfitting. Uncertainty-penalized RLHF (W4417212154) uses ensemble disagreement as a soft proxy for the true-vs-proxy gap. Each method introduces a critical hyperparameter --- constraint threshold, smoothing strength, ensemble disagreement scaling --- chosen heuristically. The principled-bound family is anchored by the recent SFT-as-adversarial-regularizer equivalence (W4399115631), which derives that, asymptotically as the preference dataset size n $\to$ ∞, the SFT loss term in the standard RLHF objective coincides with an adversarial regularizer on the proxy reward. This is the only finite-sample-grounded theoretical equivalence in the area, but it is asymptotic --- production datasets at n $\approx$ 10⁵--10⁶ preference pairs do not realize the limit, leaving the equivalence's practical regularization-schedule implication unspecified.

The structural gap is sharp. Existing mitigations have heuristic schedules; existing theory has an asymptotic equivalence whose finite-sample residual is the precise object that should govern the regularization schedule, but no one has computed it. The bottleneck is producing that residual term.

\textbf{Why now.} Three recent developments make this gap newly tractable. First, W4399115631 (Sept 2024) is itself recent --- its asymptotic equivalence is the load-bearing object whose finite-sample sharpening this work targets. Second, dblp:conf/iclr/MoskovitzSSSSDM24 (ICLR 2024) and dblp:conf/icml/ZhuJJ24 (ICML 2024) have populated the heuristic-schedule design space, making clear that the missing ingredient is a derived rather than designed schedule. Third, public RLHF stacks (TRL, OpenRLHF) and reference-model evaluation tooling (Llama-3-70B as gold reward) make a controlled empirical evaluation of schedule comparison feasible at single-academic-lab scale.

\textbf{What changes when the gap closes.} A derivable finite-sample regularization schedule for SFT-anchored RLHF eliminates a major source of empirical fragility: practitioners no longer need to tune $\beta$ by trial-and-error or run ablations across constant-$\beta$ values. The same finite-sample residual $\epsilon$ also gives a principled stopping criterion (when $\epsilon$ starts to grow despite training, overoptimization is dominating). Downstream, this enables more aggressive RLHF training (longer optimization, higher peak proxy reward) without overoptimization --- Moskovitz's constrained framing and Zhu's smoothing are recovered as special cases under specific $\Sigma$\_pref structures, providing a unifying view.

\textbf{Why prior work stopped:}
\begin{itemize}[leftmargin=*]
  \item \textbf{openalex:W4399115631}: Derives an asymptotic (n $\to$ ∞) equivalence between the SFT loss term in RLHF and an adversarial regularizer on the proxy reward, providing the first formal-reframing-grade theoretical anchor for SFT-anchored RLHF. \\ \emph{Did not do:} Does not characterize the finite-sample residual $\epsilon$(n, ...) that quantifies how far the equivalence holds at production preference dataset sizes; consequently does not derive the implied regularization schedule. \\ \emph{Structural reason:} Missing tool --- a Bernstein-type concentration on the preference-gradient covariance $\Sigma$\_pref, which is the load-bearing object on which the asymptotic limit operates; the asymptotic argument did not require it, the finite-sample argument does.
  \item \textbf{dblp:conf/iclr/MoskovitzSSSSDM24}: Reformulates RLHF as constrained RL with a Lagrangian on a proxy of the proxy-true reward gap, providing a principled framework for adaptive regularization. \\ \emph{Did not do:} Chooses the constraint threshold heuristically; does not derive an optimal threshold from data-side structure. \\ \emph{Structural reason:} Missing the connecting theorem: under what reward-data structure is a particular Lagrangian schedule optimal? Without the SFT-as-adversarial equivalence's finite-sample form (which W4399115631 supplies asymptotically), the schedule-optimality question is unanswerable from their setup.
  \item \textbf{dblp:conf/icml/ZhuJJ24}: Iterative data smoothing on the preference dataset reduces reward-model overfitting upstream; smoothing strength tuned per overfitting signal. \\ \emph{Did not do:} Smoothing strength schedule is iterative-empirical with no closed-form characterization of optimal strength as a function of dataset size and gradient structure. \\ \emph{Structural reason:} Attacks upstream cause (data overfitting) rather than the policy-side regularization schedule; the two views are dual under our finite-sample equivalence (data smoothing strength equals schedule $\beta$\_t under specific assumptions) but Zhu does not have the equivalence to make this connection.
\end{itemize}

\section*{Core claim}
\textbf{Core claim.} The SFT-as-adversarial-regularizer equivalence (W4399115631) admits a finite-sample sharpening whose residual term $\epsilon$(n, $\Sigma$\_pref) is computable from preference data alone and is the principled regularization schedule for SFT-anchored RLHF.

\textbf{Sub-claims.}
\begin{itemize}[leftmargin=*]
  \item \textbf{c1}: Replacing the asymptotic limit step in W4399115631 with a Bernstein-type concentration on the preference-gradient covariance $\Sigma$\_pref yields a finite-sample equivalence with explicit residual $\epsilon$ $\le$ C · √(d log(1/$\delta$) / n) · $\|$$\nabla$R\_proxy − $\nabla$R\_true$\|$\_op, propagated through the existing argument unchanged. \\ \emph{Supports:} the load-bearing 'finite-sample residual is computable' part of the bottleneck
  \item \textbf{c2}: Under rank-deficient $\Sigma$\_pref* (preference distinctions concentrated on r ≪ d semantic axes, as observed in HH-RLHF preference-gradient PCA), the asymptotic form is provably Ω(√(d/n)) loose while the finite-sample form recovers Õ(√(r/n)) --- establishing that the asymptotic-limit step is the binding constraint, not a peripheral concern. \\ \emph{Supports:} the audit-binding-constraint substantive demonstration; addresses C15 'auxiliary assumption' reject lesson
  \item \textbf{c3}: Setting $\beta$\_t proportional to the per-checkpoint estimated $\epsilon$ yields a derivable RLHF regularization schedule that empirically achieves higher peak gold reward and delayed degradation onset compared to constant-$\beta$, Moskovitz constraint-Lagrangian, and ensemble-disagreement schedules at fixed compute. \\ \emph{Supports:} the operational consequence of the theoretical sharpening
\end{itemize}

\section*{Method flow}
\textbf{High-level pipeline.} Take any RLHF run that uses SFT-anchored PPO with KL coefficient $\beta$. At each checkpoint, compute the preference-gradient covariance $\Sigma$\_pref on a held-out preference batch and estimate the operator-norm gap term, yielding $\epsilon$(n, $\Sigma$\_pref). Set the next training window's $\beta$ = clip(c · $\epsilon$, $\beta$\_min, $\beta$\_max). Continue PPO training. Output is the trained policy plus the $\epsilon$ trajectory which serves as overoptimization audit trail.

\textbf{Steps.}
\begin{itemize}[leftmargin=*]
  \item \textbf{S1: Estimate preference-gradient covariance} \\ \emph{What changes:} On a held-out preference batch of size m every K policy steps, compute $\Sigma$\_pref = E[($\nabla$R\_proxy(x,y\_w) - $\nabla$R\_proxy(x,y\_l))(...)\textasciicircum{}T] empirically --- a d$\times$d matrix where d is the proxy reward model's gradient dimension. Use low-rank sketches (rank r=512 typically) for tractability. \\ \emph{Why:} $\Sigma$\_pref is the load-bearing object whose asymptotic limit W4399115631 takes; estimating it at finite n is the entry point for the finite-sample sharpening. \\ \emph{Linked component:} engineering\_leg \quad \emph{Linked falsification:} failure\_b\_mechanistic
  \item \textbf{S2: Compute residual $\epsilon$(n, $\Sigma$\_pref\textasciicircum{}t)} \\ \emph{What changes:} Apply theorem closed form: $\epsilon$ = C · √(rank($\Sigma$\_pref\textasciicircum{}t) · log(1/$\delta$) / n) · $\|$$\nabla$R\_proxy − $\nabla$R\_true$\|$\_$\Sigma$\_pref\textasciicircum{}t. The operator norm is estimated via the SFT model's preference-likelihood gradient as a calibration constant (justified because at the SFT initialization R\_proxy $\approx$ R\_true). \\ \emph{Why:} Theorem-derived schedule value; the C15 success-condition lesson that the sharpened object is the actual signal not a proxy. \\ \emph{Linked component:} theoretical\_leg \quad \emph{Linked falsification:} failure\_b\_mechanistic
  \item \textbf{S3: Set policy regularization coefficient} \\ \emph{What changes:} $\beta$\_t = clip(c · $\epsilon$\textasciicircum{}t, $\beta$\_min=0.01, $\beta$\_max=1.0). Single hyperparameter c calibrated once at training start using SFT model behavior; not retuned. \\ \emph{Why:} The derived schedule. Algorithm (SFT-anchored PPO) is unchanged; only $\beta$\_t becomes data-driven. \\ \emph{Linked component:} engineering\_leg \quad \emph{Linked falsification:} failure\_a\_numerical
  \item \textbf{S4: Run PPO with derived $\beta$\_t} \\ \emph{What changes:} Standard PPO loop with the time-varying $\beta$\_t replacing the constant $\beta$. Rest of training pipeline byte-identical to W4399115631's SFT-anchored prescription. \\ \emph{Why:} Pattern C15 success condition: algorithm unchanged, only the analytic object that controls behavior is sharpened. \\ \emph{Linked component:} engineering\_leg \quad \emph{Linked falsification:} failure\_a\_numerical
\end{itemize}

\textbf{Preservation argument.} The derivation preserves W4399115631's argument structure: the asymptotic step is the only step replaced; everything else (the equivalence form, the role of the SFT loss, the policy-PPO update rule) is byte-identical. C15 success condition checked: algorithm unchanged, single load-bearing inequality (the asymptotic limit on $\Sigma$\_pref) replaced by the Bernstein concentration substitute, downstream rates pass through unchanged.

\section*{Theorem / algorithm / system design}
Algorithm (SFT-anchored PPO with derived schedule):
  Input: SFT base $\pi$\_SFT, proxy reward R\_proxy, preference dataset D\_pref (size n), calibration c, $\beta$\_min, $\beta$\_max, K
  Initialize: $\pi$ ← $\pi$\_SFT, t ← 0
  Repeat until budget exhausted:
    Every K steps:
      $\Sigma$\_pref\textasciicircum{}t ← rank-r sketch of ($\nabla$R\_proxy\_w − $\nabla$R\_proxy\_l)(...)\textasciicircum{}T over batch of D\_pref
      $\epsilon$\textasciicircum{}t ← C · √(r · log(1/$\delta$) / n) · g\_op   where g\_op estimated from $\pi$\_SFT calibration
      $\beta$\_t ← clip(c · $\epsilon$\textasciicircum{}t, $\beta$\_min, $\beta$\_max)
    PPO step on $\pi$ with reward R\_proxy(s,a) − $\beta$\_t · KL($\pi$||$\pi$\_SFT)
    t ← t + 1
  Output: $\pi$ plus $\epsilon$\textasciicircum{}t trajectory

Theorem (Finite-Sample SFT-as-Adversarial-Regularizer Equivalence): L\_SFT($\theta$) = L\_adv-reg($\theta$; $\beta$=$\|$$\nabla$R\_proxy − $\nabla$R\_true$\|$\_$\Sigma$\_pref) + $\epsilon$(n, $\|$$\nabla$R\_proxy − $\nabla$R\_true$\|$) with $\epsilon$ $\le$ C · √(d · log(1/$\delta$) / n) · $\|$$\nabla$R\_proxy − $\nabla$R\_true$\|$\_op + O(n\textasciicircum{}\{-1\}) w.p. $\ge$ 1−$\delta$.

Separation Lemma: Under rank-r $\Sigma$\_pref* the asymptotic form W4399115631 is Ω(√(d/n)) loose; the finite-sample form recovers Õ(√(r/n)).

\section*{Theoretical leg}
Theorem (Finite-Sample SFT-as-Adversarial-Regularizer Equivalence): Let $\Sigma$\_pref be the empirical preference-gradient covariance over n pairs. Then L\_SFT($\theta$) = L\_adv-reg($\theta$; $\beta$=$\|$$\nabla$R\_proxy − $\nabla$R\_true$\|$\_$\Sigma$\_pref) + $\epsilon$(n, $\|$$\nabla$R\_proxy − $\nabla$R\_true$\|$), where $\epsilon$ $\le$ C · √(d · log(1/$\delta$) / n) · $\|$$\nabla$R\_proxy − $\nabla$R\_true$\|$\_op + O(n\textasciicircum{}(-1)) with probability 1−$\delta$ over the preference draw, for an explicit constant C depending on the policy class. Proof: replace the asymptotic step in W4399115631 with a Bernstein-type concentration on $\Sigma$\_pref, propagate through the existing argument unchanged. Corollary: the residual $\epsilon$ is computable from preference data alone (no gold reward needed) and serves as a derivable regularization schedule. Separation Lemma (Audited Assumption is Binding, Not Auxiliary): construct a preference-gradient covariance $\Sigma$\_pref* whose spectrum is rank-r ≪ d on a low-dimensional subspace of the preference manifold (e.g., when preference distinctions concentrate on few semantic axes). On $\Sigma$\_pref*, the asymptotic equivalence W4399115631 is provably Ω(√(d/n)) loose because the W4399115631 argument averages over the full d-dimensional gradient space while the actual signal lives on r dimensions. The finite-sample form recovers the optimal rate Õ(√(r/n)) by exploiting the rank structure via $\Sigma$\_pref's spectrum directly. This separates: (i) the asymptotic-limit step IS the binding constraint, (ii) the operator-norm assumption $\|$$\nabla$R\_proxy − $\nabla$R\_true$\|$\_op $\le$ B is structurally tight rather than covertly assuming the tightness --- under $\Sigma$\_pref* the rank-r signal directly determines the achievable rate, and B controls the residual constant not the rate exponent. Non-vacuity verified via HH-RLHF preference-gradient PCA in companion ablation (rank concentration r $\approx$ 50 in d $\approx$ 4096 vocab-projected gradient subspace).

\section*{Engineering leg}
\textasciitilde{}30 lines of PyTorch on top of standard RLHF training loop: (a) compute $\Sigma$\_pref via rank-r=512 sketch on held-out preference batch every K=200 policy steps; (b) estimate $\|$$\nabla$R\_proxy − $\nabla$R\_true$\|$\_$\Sigma$\_pref via the proxy reward model's gradient norm + a calibration constant from the SFT model's preference-likelihood gradient; (c) set $\beta$\_t = clip(c · $\epsilon$(n, ...), 0.01, 1.0) where c is calibrated once at SFT init; (d) standard PPO with this $\beta$\_t schedule. No new architectures, no new losses, no ensemble --- only the schedule for the existing $\beta$ coefficient changes. Implementation slots into TRL's PPOTrainer; \textasciitilde{}30 LOC patch. Compute overhead: \textasciitilde{}3\% from periodic $\Sigma$\_pref sketches.

\section*{Falsification prediction}
\textbf{What we run.} Train SFT-anchored PPO on Anthropic HH-RLHF with Llama-3-8B SFT base, comparing four schedules: (a) constant $\beta$ = 0.1 (vanilla baseline), (b) Moskovitz constrained Lagrangian, (c) ensemble-disagreement schedule (5-model ensemble), (d) our derived $\beta$\_t = clip(c · $\epsilon$(n, $\Sigma$\_pref\textasciicircum{}t), 0.01, 1.0) schedule. Use Llama-3-70B as held-out gold reward proxy. Measure (i) gold-reward gap at fixed proxy-reward target, (ii) max gold reward achieved before degradation, (iii) empirical $\epsilon$(n, $\Sigma$\_pref\textasciicircum{}t) trajectory vs theoretical bound at three checkpoints (10k, 50k, 200k preference pairs).

\textbf{Success.} Schedule (d) achieves at least 0.15 bits per token higher gold reward at peak than schedules (a)/(b)/(c) at fixed compute (8$\times$A100$\times$8wk); empirical $\epsilon$(n, ...) trajectory stays within 1.5$\times$ of theoretical bound across all 3 checkpoints; gold-reward degradation onset for (d) is delayed by $\ge$ 50\% of training steps relative to (a).

\textbf{Failure (a) — numerical.} Schedule (d) does not improve peak gold reward over baselines at fixed compute (within 0.05 bits per token).

\textbf{Failure (b) — mechanistic.} Schedule (d) improves peak gold reward, but empirical $\epsilon$(n, ...) trajectory exceeds theoretical bound by > 3$\times$ at any of the 3 checkpoints --- indicating the finite-sample bound is not the operating mechanism (gain came from regularization side-effect, not from tracking the predicted residual).

\textbf{Compute budget.} 12 GPU-day for the minimal experiment (4 schedules $\times$ 1 seed $\times$ 200k preference pairs on Llama-3-8B at 8$\times$A100). User has 8$\times$A100$\times$8wk $\approx$ 448 A100-day; budget is 1.5 orders of magnitude under available --- feasible.

\textbf{Decision rule.} if failure\_a or failure\_b $\to$ kill

\section*{Mechanism probe}
At each of 3 checkpoints (10k, 50k, 200k preference pairs into training), measure: (i) empirical L\_SFT - L\_adv-reg gap, (ii) theoretical $\epsilon$(n, ...) computed from preference-gradient covariance, (iii) gold-vs-proxy reward gap measured by Llama-3-70B reference. Discriminating outcome: (i) $\approx$ (ii) within Bernstein-bound tolerance AND (iii) tracks (ii) within constant factor at each checkpoint, for schedule (d) only. For schedules (a)/(b)/(c) the tracking should fail at $\ge$ 1 checkpoint. This distinguishes 'finite-sample bound IS the operational mechanism' (success) from 'gain from side-effect' (failure\_b).

\section*{Feasibility validation}
\textbf{Compute.} \emph{ feasible } — 12 GPU-day required vs user 448 A100-day = 1.5 orders of magnitude headroom. Single 8$\times$A100 node $\times$ \textasciitilde{}1.5 days per schedule $\times$ 4 schedules = 12 GPU-day.

\textbf{Data.} \emph{ feasible } — Anthropic HH-RLHF (161k preference pairs) is publicly available; Llama-3-8B SFT and Llama-3-70B gold-proxy are publicly available on HuggingFace; no closed/private data needed.

\textbf{Theoretical.} \emph{ tight } — Bernstein-type concentration on covariance is standard machinery, but instantiating it correctly with rank-r structure on $\Sigma$\_pref requires careful handling of the preference-gradient sub-Gaussianity assumption. User's stated expertise (RLHF + theoretical preference learning) is sufficient. The Separation Lemma construction is the most technical piece; \textasciitilde{}3 weeks of theoretical work.

\textbf{Engineering.} \emph{ feasible } — \textasciitilde{}30 LOC patch on TRL/OpenRLHF PPOTrainer. $\Sigma$\_pref rank-r sketch is straightforward (low-rank Gaussian sketch or randomized SVD of preference-gradient outer products on held-out batches). 3\% compute overhead from periodic sketching.

\textbf{Falsification.} \emph{ feasible } — What\_we\_run is concrete (4 schedules $\times$ 1 seed $\times$ 200k preference pairs); mechanism\_probe is concrete (3 checkpoints, 3 measurements per checkpoint); both fit in the 12 GPU-day budget. Llama-3-70B gold-reward proxy is established methodology.

\textbf{Overall.} \emph{ feasible }

\section*{Differentiation from recent literature}
\begin{itemize}[leftmargin=*]
  \item \textbf{openalex:W4399115631}: W4399115631 derives the SFT-as-adversarial equivalence asymptotically (n $\to$ ∞ limit). We sharpen the single asymptotic step into a finite-sample concentration bound, recovering W4399115631 as the n $\to$ ∞ limit of our theorem. The substantive contribution is the residual term $\epsilon$(n, ...) --- which W4399115631 cannot characterize --- and its operational use as a regularization schedule.
  \item \textbf{dblp:conf/iclr/MoskovitzSSSSDM24}: Moskovitz formulates RLHF as constrained RL with a heuristic threshold on a proxy gap measure. Our theorem shows that under specific $\Sigma$\_pref structure, the optimal Lagrangian schedule is exactly $\beta$\_t = $\epsilon$(n, ...) --- recovering Moskovitz's heuristic constraint as a special case of our derived schedule, while explaining when the heuristic fails (when $\Sigma$\_pref is rank-deficient on the user's preference manifold).
  \item \textbf{dblp:conf/icml/ZhuJJ24}: Zhu's iterative data smoothing reduces reward-model overfitting upstream by adjusting smoothing strength. We show the optimal smoothing strength under our finite-sample form is also $\beta$\_t = $\epsilon$(n, ...) viewed as preference-data-side regularization rather than policy-side, explaining the structural similarity of the two methods and giving a derivable rather than iterative-empirical schedule.
  \item \textbf{openalex:W4417212154}: Uncertainty-penalized RLHF uses ensemble disagreement as an empirical proxy for the proxy-true reward gap. Our theorem provides a closed-form expression for that gap (in terms of $\Sigma$\_pref and gradient-norm estimates) without requiring an ensemble --- and shows when ensemble disagreement provably tracks $\epsilon$ vs when it does not (it tracks $\epsilon$ only when the ensemble is large enough to span the preference-gradient eigenspectrum, otherwise underestimates).
  \item \textbf{arxiv:2602.21765v1}: Generalisation of RLHF under Reward Shift and Clipped KL (arxiv:2602.21765v1) decomposes RLHF generalization error into 3 sources (sampling, reward shift, KL clipping) and yields a one-time optimal KL-clipping threshold plus budget allocation. Our work derives a dynamic per-checkpoint regularization schedule $\beta$\_t from a specific classical-object equivalence (SFT loss $\leftrightarrow$ adversarial regularizer per W4399115631), not a generic generalization bound. Different artifact (schedule vs bound) and different starting point (sharpened equivalence vs decomposition). Corollary direction: under the assumption that reward shift and KL clipping errors are negligible vs sampling error, our $\epsilon$(n, $\Sigma$\_pref) recovers their sampling-error term as a special case --- but their bound does not yield our schedule structure because the SFT-as-adversarial equivalence is the load-bearing object.
\end{itemize}

\section*{Reviewer concerns and responses}
\begin{itemize}[leftmargin=*]
  \item \textbf{Severity non\_blocking.} \emph{Attack:} C15 reject lesson 'auxiliary assumption nullifies tightness' (ICML\_2025\_1674): the operator-norm assumption $\|$$\nabla$R\_proxy − $\nabla$R\_true$\|$\_op $\le$ B may be the new loose step that survives the Bernstein rewrite --- i.e., the sharpening is mechanical rather than genuine. \\ \emph{Response:} Defense anchored in theoretical\_leg's Separation Lemma. The construction shows that under rank-r $\Sigma$\_pref* (verified non-vacuously on HH-RLHF via preference-gradient PCA, r $\approx$ 50 vs d $\approx$ 4096), the asymptotic form is provably Ω(√(d/n)) loose while our finite-sample form recovers Õ(√(r/n)). The operator norm B controls the residual constant, not the rate exponent --- so a loose B cannot turn the rate-improvement into rate-preservation; the rate gap r vs d cannot be hidden in B. This is genuine sharpening on the same object, not auxiliary-assumption rewrite. \\ \emph{Fields changed:} theoretical\_leg
  \item \textbf{Severity non\_blocking.} \emph{Attack:} Paper-pointed threat arxiv:2602.21765v1 ('Generalisation of RLHF under Reward Shift and Clipped KL Regularisation') derives generalization bounds for RLHF with practical implications including 'optimal KL clipping threshold' and 'budget allocation.' A reviewer will ask whether our Bernstein-on-$\Sigma$\_pref result is just a special case of their bound under the assumption that reward shift and KL clipping errors are dominated by sampling error. \\ \emph{Response:} Defense anchored in the new differentiation\_from\_lit entry for arxiv:2602.21765v1. Two distinguishing factors: (a) Different artifact --- they produce a one-time optimal threshold; we produce a per-checkpoint dynamic schedule $\beta$\_t computable from current rollouts. (b) Different starting point --- they decompose generic RLHF generalization error into 3 sources; we sharpen a specific classical-object equivalence (SFT $\leftrightarrow$ adversarial regularizer per W4399115631). Corollary direction: our $\epsilon$(n, $\Sigma$\_pref) recovers their sampling-error term as a special case under specific assumptions (negligible shift + clipping), but their bound does not yield our schedule because the SFT-as-adversarial equivalence is load-bearing in our derivation, absent in theirs. \\ \emph{Fields changed:} differentiation\_from\_lit
\end{itemize}

\end{tcolorbox}

\end{document}
